% =====================================================================
% config/commands.tex
% PT: Comandos auxiliares próprios. Acrescente aqui os seus comandos
%     (\newcommand) para os manter num só sítio.
% EN: Helper commands of your own. Add your commands (\newcommand) here
%     to keep them in one place.
% =====================================================================

% PT: Caixa vazia para marcar uma figura que ainda não existe.
%     Uso: \placeholderfigure{<largura>}{<altura>}
%     Ex.: \placeholderfigure{0.8\textwidth}{5cm}
% EN: Empty box to mark a figure that does not exist yet.
%     Usage: \placeholderfigure{<width>}{<height>}
%     E.g.: \placeholderfigure{0.8\textwidth}{5cm}
\newcommand{\placeholderfigure}[2]{%
  \begin{center}
    \fbox{\parbox[c][#2][c]{#1}{\centering Figure placeholder / Figura provisória\\[0.5em]
    Replace with the final image. / Substituir pela imagem final.}}
  \end{center}
}

% PT: Nota de trabalho a azul, para lembretes durante a escrita.
%     Apague todas antes da entrega. Uso: \todosection{texto}
% EN: Work note in blue, for reminders while writing.
%     Remove them all before submission. Usage: \todosection{text}
\newcommand{\todosection}[1]{%
  \begin{quote}
    \color{blue!70!black}\textbf{TODO:} #1
  \end{quote}
}

% PT: Questões de investigação e objetivos numerados.
%     Uso: \researchq{1}{Texto da questão}  ->  RQ1. Texto da questão
% EN: Numbered research questions and objectives.
%     Usage: \researchq{1}{Question text}  ->  RQ1. Question text
\newcommand{\researchq}[2]{\textbf{RQ#1.} #2}
\newcommand{\researchobj}[2]{\textbf{OBJ#1.} #2}

% PT: Coluna Y para tabularx: ocupa a largura disponível e alinha à esquerda.
%     Uso: \begin{tabularx}{\textwidth}{lY} ... \end{tabularx}
% EN: Y column for tabularx: fills the available width, left-aligned.
%     Usage: \begin{tabularx}{\textwidth}{lY} ... \end{tabularx}
\newcolumntype{Y}{>{\raggedright\arraybackslash}X}
